\subsection{What Limits the Accuracy}
\label{subsec:disc-accuracy}
The experiments separate two regimes that the bound of Section~\ref{sec:bounds} distinguishes. Measurement noise, at an SNR of several thousand per record, would allow millimeter accuracy; the centimeter errors measured on D2 and D3 are instead set by systematic errors of the model. The chain of Table~\ref{tab:ablation} identifies the main ones, and every step corresponds to a physical element: the angular response of the receiver, the angular zero of the stage, and the position and pattern of the LED. The remaining error has a clear signature, a vertical bias that changes with the distance to the LED and a horizontal error that repeats at the same horizontal positions, and the calibration study shows that the single-order Lambertian pattern alone cannot go below 3.15~cm. A pattern model with an angle-dependent order, a spline, or a model including the tilt of the LED, calibrated on two planes at different heights so that the LED height also becomes identifiable, is the natural next step.

\subsection{Calibration Effort}
The calibration required here is modest and, except for two tape measurements, does not use the positioning data. The LED model needs a few flat-PD measurements spread over the working area: nine points on a regular grid suffice, provided they cover the emission angles used for positioning. The angular zero of the stage is a property of the receiver assembly; it was calibrated with the shape test at known positions, which a manufacturer or a one-time setup procedure can provide, and its observed drift ($\approx0.5^\circ$ over ten hours) suggests that an absolute angular reference on the stage would be preferable to an incremental one. The receiver response was taken from the literature; the shape test confirmed it, but a characterization of the individual device would remove one more assumption.

\subsection{Reliability of the Orientation}
The method trusts the orientations reported by the actuator, and Section~\ref{subsec:res-orient} quantifies how much: about $0.2^\circ$ keeps the orientation errors small against the model bias, and errors common to all orientations, such as a drifting angular zero, are the most harmful. Two faults of the rotation stage occurred during this work, a tilt axis that moved by itself and a rotation that slipped while the encoder kept reporting the commanded angles. Both produced errors several times larger than nominal, and both were detectable in the residual of the fit, without ground truth. A practical system should monitor this residual and, when it rises persistently, stop trusting the orientation set. The same reasoning applies if the stage is replaced by an inertial measurement unit, whose errors enter the bound through~\eqref{eq:fim-orientation}.

\subsection{\texorpdfstring{\add{Towards Mobile Receivers}}{Towards Mobile Receivers}}
\label{subsec:disc-mobile}
\add{The main motivation for a single reoriented PD is its minimal hardware. A platform that turns about its vertical axis, such as a drone in hover, realizes the cone of orientations by yawing, with the normals given by its IMU: one tilted PD, mounted close to the yaw axis so that its optical center barely moves, then provides 3D positioning at a small cost in size, weight and power, which are the scarce resources of such platforms. Since 10~ms of signal per orientation suffice (Section~\ref{subsec:res-window}) and the orientations need not be uniform (Section~\ref{subsec:res-K}), the acquisition is limited only by the rotation of the platform.} \add{The normals then come from the attitude of the platform, expressed in the global frame of the room (Section~\ref{subsec:geometry}), as the angular zero of the stage was here: roll and pitch follow from gravity, and the heading from an initial alignment with the axes of the room, which the shape test~\eqref{eq:shape} can provide at a known position. A residual heading error only rotates the estimate about the vertical through the LED, by 1.7~cm per degree at 1~m from it (1.2~cm on average over D2), without changing the distance or the height. Several LEDs make the heading observable from their directions; this and other heading references will be explored in future work.}

\subsection{Limitations}
The ground truth relies on the gantry, on tape measurements of the LED plane and of the carriage height, and on the CAD elevation of the PD center; the horizontal position of each LED is inferred from the calibration rather than surveyed. The experiments cover one room, one LED height and a $1.2\times1.2\times0.36$~m volume below the positioning LED; larger volumes, other heights and other receivers remain to be tested, and the results of Section~\ref{subsec:res-leds} show that a tilted LED must be modelled or realigned. \add{Reflections from the walls and the furniture were not modelled; whatever they contribute is part of the model bias of Section~\ref{subsec:disc-accuracy}, and a channel model with first-order reflections is left for future work.} The acquisition time, about four minutes per position, is dictated by the rotation stage; a faster actuator\del{,}\add{ or} continuous rotation with on-the-fly sampling\del{, or a receiver with several fixed PDs} would remove this limitation and open the way to moving receivers.\add{ The model, the bounds and the estimators require only $K$ known normals, so they also apply to receivers that realize them simultaneously with several fixed PDs, at the price of matching the gains and responses of the detectors.} Finally, only the single-LED configuration was evaluated; when several LEDs are visible, their directions alone could be intersected without any gain or pattern calibration, which is left for future work.
